\documentclass[10pt, aspectratio=169]{beamer}
% Preámbulo modular para presentaciones Beamer (Modelación Estadística)
% Diseño y paleta institucional del Tecnológico de Monterrey

\documentclass[aspectratio=169,xcolor={dvipsnames,table}]{beamer}

\usepackage[utf8]{inputenc}
\usepackage[spanish,mexico]{babel}
\usepackage{amsmath,amssymb,amsfonts,mathtools}
\usepackage{graphicx}
\usepackage{listings}
\usepackage{booktabs}
\usepackage{tikz}
\usetikzlibrary{matrix,arrows,decorations.pathmorphing,shapes.geometric,positioning}

% Paleta Institucional Tecnológico de Monterrey
\definecolor{TecAzul}{HTML}{0039A6}       % Azul TEC: color primario institucional
\definecolor{TecAzulOscuro}{HTML}{1A2E51} % Azul marino oscuro
\definecolor{TecRojo}{HTML}{EC2661}       % Rojo de acento (paleta miTec)
\definecolor{TecGris}{HTML}{646464}       % Gris neutro para comentarios y subtítulos
\definecolor{TecFondoBloque}{HTML}{F4F6F9}% Fondo suave para bloques

% Configuración de Tema Beamer
\usetheme{Boadilla}
\usecolortheme[named=TecAzulOscuro]{structure}
\setbeamercolor{alerted text}{fg=TecRojo}
\setbeamercolor{palette primary}{bg=TecAzulOscuro,fg=white}
\setbeamercolor{palette secondary}{bg=TecAzul,fg=white}
\setbeamercolor{palette tertiary}{bg=TecAzulOscuro!80!black,fg=white}
\setbeamercolor{title}{fg=TecAzulOscuro,bg=TecFondoBloque}
\setbeamercolor{frametitle}{fg=TecAzulOscuro,bg=TecFondoBloque}

% Bloques personalizados elegantes
\setbeamercolor{block title}{bg=TecAzulOscuro,fg=white}
\setbeamercolor{block body}{bg=TecFondoBloque,fg=black}
\setbeamercolor{block title alerted}{bg=TecRojo,fg=white}
\setbeamercolor{block body alerted}{bg=TecRojo!8,fg=black}
\setbeamercolor{block title example}{bg=TecAzul,fg=white}
\setbeamercolor{block body example}{bg=TecAzul!8,fg=black}

% Redondear bordes y añadir sombra sutil
\setbeamertemplate{blocks}[rounded][shadow=true]
\setbeamertemplate{navigation symbols}{}

% Pie de página limpio con número de diapositiva
\setbeamertemplate{footline}{
  \leavevmode%
  \hbox{%
  \begin{beamercolorbox}[wd=.7\paperwidth,ht=2.25ex,dp=1ex,left]{author in head/foot}%
    \hspace*{2ex}\insertshortauthor~~(\insertshortinstitute) \hfill \insertshorttitle
  \end{beamercolorbox}%
  \begin{beamercolorbox}[wd=.3\paperwidth,ht=2.25ex,dp=1ex,right]{date in head/foot}%
    \insertshortdate{}\hspace*{2em}
    \insertframenumber{} / \inserttotalframenumber\hspace*{2ex}
  \end{beamercolorbox}}%
  \vskip0pt%
}

% Configuración de Listings de Python para visualización pedagógica de código
\lstset{
    language=Python,
    basicstyle=\ttfamily\footnotesize,
    keywordstyle=\color{TecAzulOscuro}\bfseries,
    stringstyle=\color{TecRojo},
    commentstyle=\color{TecGris}\itshape,
    showstringspaces=false,
    breaklines=true,
    frame=lines,
    rulecolor=\color{TecAzulOscuro!30},
    backgroundcolor=\color{TecFondoBloque},
    aboveskip=0.5em,
    belowskip=0.5em,
    literate={á}{{\'a}}1 {é}{{\'e}}1 {í}{{\'i}}1 {ó}{{\'o}}1 {ú}{{\'u}}1
             {Á}{{\'A}}1 {É}{{\'E}}1 {Í}{{\'I}}1 {Ó}{{\'O}}1 {Ú}{{\'U}}1
             {ñ}{{\~n}}1 {Ñ}{{\~N}}1 {¿}{{?`}}1 {¡}{{!`}}1
}

% Metadatos institucionales por defecto
\title[Modelación Estadística]{Modelación Estadística para Ciencia de Datos e Ingeniería}
\author[J. Castillo Colmenares]{Juliho David Castillo Colmenares}
\institute[Tec de Monterrey]{Tecnológico de Monterrey \\ Escuela de Ingeniería y Ciencias}
\date{\today}

% Comandos y macros estadísticas compartidas para las presentaciones Beamer (Metropolis)

% Conjuntos numéricos básicos
\newcommand{\R}{\mathbb{R}}
\newcommand{\N}{\mathbb{N}}
\newcommand{\Q}{\mathbb{Q}}
\newcommand{\Z}{\mathbb{Z}}
\newcommand{\set}[1]{\left\{ #1 \right\}}
\newcommand{\abs}[1]{\left|#1\right|}
\newcommand{\norm}[1]{\left\| #1 \right\|}

% Comandos estadísticos y probabilísticos del curso
\newcommand{\Var}{\operatorname{Var}}
\newcommand{\cov}{\operatorname{Cov}}
\newcommand{\corr}{\rho}
\newcommand{\comb}[2]{\begin{pmatrix} #1 \\ #2 \end{pmatrix}}
\newcommand{\s}{\sigma}
\newcommand{\particion}{\mathfrak{P}}
\newcommand{\card}[1]{\#\left( #1 \right)}
\newcommand{\seq}[3]{\set{#1}_{#2}^{#3}}
\newcommand{\E}{\mathbb{E}}
\newcommand{\Prob}{\mathbb{P}}

% Entornos pedagógicos y cajas en español académico natural
\newenvironment{discusion}{%
  \begin{alertblock}{Pregunta de Discusión}%
}{%
  \end{alertblock}%
}

\newenvironment{reflexion}{%
  \begin{alertblock}{Reflexión para el Aula}%
}{%
  \end{alertblock}%
}

\newenvironment{paradoja}{%
  \begin{alertblock}{Paradoja de Exploración}%
}{%
  \end{alertblock}%
}

\newenvironment{motivacion}{%
  \begin{exampleblock}{Motivación: Ciencia de Datos en la Práctica}%
}{%
  \end{exampleblock}%
}

\newenvironment{retoclase}{%
  \begin{block}{Reto Rápido en Equipo}%
}{%
  \end{block}%
}


\title[$F$ de Fisher-Snedecor]{Sección 05.05: Distribución $F$ de Fisher-Snedecor}
\subtitle{Comparación de Varianzas y Análisis de Varianza (ANOVA)}
\author[J. Castillo Colmenares]{Juliho Castillo Colmenares}
\institute{Tecnológico de Monterrey}
\date{\vspace{-1.2cm}}

\begin{document}

% Slide 1: Portada
\begin{frame}[plain]
  \titlepage
\end{frame}

% Slide 2: Hoja de Ruta
\begin{frame}{Hoja de Ruta --- Capítulo 05: Distribuciones de Muestreo}
  \footnotesize
  ¿Dónde nos encontramos en el desarrollo modular del capítulo? \pause
  \begin{itemize}\itemsep=0.08em
    \item \textbf{05.01} Muestreo Aleatorio Simple, Media y Varianza Muestral Insesgada \pause
    \item \textbf{05.02} Teorema del Límite Central Asintótico \pause
    \item \textbf{05.03} Distribución Chi-Cuadrada y Varianza Muestral \pause
    \item \textbf{05.04} Distribución $t$ de Student y Muestras Pequeñas \pause
    \item \textbf{05.05} Distribución $F$ de Fisher-Snedecor \emph{(Hoy --- Cierre del Capítulo)}
  \end{itemize}
  \vspace{-0.05cm}
  \begin{block}{Objetivo de la Sesión}
    Formalizar la distribución $F$ como cociente de chi-cuadradas independientes, dominar la prueba $F$ de igualdad de varianzas y su intervalo de confianza asociado, y aplicar el estadístico $F$ al Análisis de Varianza (ANOVA) para comparar múltiples medias simultáneamente.
  \end{block}
\end{frame}

% Slide 3: Motivación
\begin{frame}{El Tercer Miembro de la Familia $\chi^2$--$t$--$F$}
  \footnotesize
  \begin{motivacion}
    Ya vimos que $\chi^2$ caracteriza la distribución de $S^2$, y que $t$ surge al estimar $\mu$ con $\sigma$ desconocida. Hoy cerramos el trío: la distribución $F$ compara \emph{dos} varianzas muestrales, y generaliza esta idea para comparar \emph{múltiples} medias a la vez mediante ANOVA.
  \end{motivacion}

  \vspace{0.15cm}
  \pause
  \begin{itemize}\itemsep=0.1cm
    \item \textbf{Origen:} cociente de dos chi-cuadradas independientes, cada una normalizada por sus grados de libertad. \pause
    \item \textbf{Aplicación 1:} ¿son iguales las varianzas de dos poblaciones? \pause
    \item \textbf{Aplicación 2:} ¿son iguales las medias de $k\ge 3$ grupos? (ANOVA)
  \end{itemize}
\end{frame}

% Slide 4: Definición y Propiedades
\begin{frame}{Definición y Propiedades de la Distribución $F$}
  \footnotesize
  Sean $\chi^2_{d_1}$ y $\chi^2_{d_2}$ independientes. \pause

  \vspace{0.1cm}
  \begin{block}{Definición}
    $$F = \dfrac{\chi^2_{d_1}/d_1}{\chi^2_{d_2}/d_2} \sim F_{d_1,d_2}$$
  \end{block}

  \vspace{0.1cm}
  \pause
  \begin{alertblock}{Propiedades}
    \begin{itemize}\itemsep=0.05cm
      \item $F\ge 0$, asimétrica (sesgada a la derecha); $E(F)=d_2/(d_2-2)$ para $d_2>2$.
      \item Recíproco: si $F\sim F_{d_1,d_2}$, entonces $1/F\sim F_{d_2,d_1}$.
      \item \textbf{Conexión con $t$:} si $T\sim t_\nu$, entonces $T^2\sim F_{1,\nu}$.
    \end{itemize}
  \end{alertblock}
\end{frame}

% Slide 5: Prueba F para Varianzas
\begin{frame}{Prueba $F$ de Igualdad de Varianzas}
  \footnotesize
  Para dos muestras normales independientes: \pause

  \vspace{0.1cm}
  \begin{block}{Estadístico de Prueba}
    $$F = \dfrac{S_1^2}{S_2^2} \sim F_{n_1-1,\,n_2-1} \quad \text{bajo } H_0:\sigma_1^2=\sigma_2^2$$
  \end{block}

  \vspace{0.1cm}
  \pause
  \begin{alertblock}{Intervalo de Confianza para $\sigma_1^2/\sigma_2^2$}
    $$\left[\dfrac{S_1^2}{S_2^2}\cdot\dfrac{1}{F_{n_1-1,n_2-1,\alpha/2}},\;\; \dfrac{S_1^2}{S_2^2}\cdot F_{n_2-1,n_1-1,\alpha/2}\right]$$
    Si el intervalo contiene a $1$, es consistente con varianzas iguales.
  \end{alertblock}
\end{frame}

% Slide 6: ANOVA
\begin{frame}{Análisis de Varianza (ANOVA)}
  \footnotesize
  ANOVA descompone la variabilidad total en componentes de tratamiento y error: \pause

  \vspace{0.1cm}
  \begin{block}{Estadístico $F$ de ANOVA}
    $$\text{SC}_{\text{trat}} = \sum_i n_i(\bar X_i-\bar X)^2, \quad \text{SC}_{\text{error}} = \sum_i\sum_j(X_{ij}-\bar X_i)^2$$
    $$F = \dfrac{\text{SC}_{\text{trat}}/(k-1)}{\text{SC}_{\text{error}}/(n-k)} \sim F_{k-1,\,n-k} \quad \text{bajo } H_0:\mu_1=\cdots=\mu_k$$
  \end{block}

  \vspace{0.1cm}
  \pause
  \begin{alertblock}{Intuición}
    Si las medias de los grupos son muy distintas entre sí (SC$_{\text{trat}}$ grande) en relación con la variabilidad dentro de cada grupo (SC$_{\text{error}}$), $F$ será grande y rechazaremos $H_0$.
  \end{alertblock}
\end{frame}

% Slide 7: Aplicaciones
\begin{frame}{Aplicaciones en Ciencia de Datos e Ingeniería}
  \footnotesize
  La distribución $F$ es la base de comparaciones multi-grupo: \pause

  \vspace{0.15cm}
  \begin{itemize}\itemsep=0.12cm
    \item \textbf{Control de calidad:} comparar la variabilidad de dos procesos de manufactura. \pause
    \item \textbf{Experimentos A/B/n:} ANOVA generaliza la prueba $t$ a más de dos variantes simultáneas. \pause
    \item \textbf{Regresión lineal:} la prueba $F$ global evalúa si el modelo en conjunto es significativo. \pause
    \item \textbf{Diseño de experimentos:} base del Capítulo 08 (ANOVA de uno y varios factores).
  \end{itemize}
\end{frame}

% Slide 8: Lab Python (Bloque 1)
\begin{frame}[fragile]{Laboratorio en Python: Propiedades de la $F$}
  \scriptsize
  Verificación de $E(F)$, la propiedad recíproca, y la identidad $T^2\sim F_{1,\nu}$ (\texttt{05.05\_fisher\_f\_distribution.py}):
  {\lstset{basicstyle=\fontsize{4pt}{4.8pt}\selectfont\ttfamily,aboveskip=0.1em,belowskip=0.1em}
  \lstinputlisting[firstline=17, lastline=36]{../../code/05_distribuciones_muestreo/05.05_fisher_f_distribution.py}}
\end{frame}

% Slide 9: Lab Python (Bloque 2)
\begin{frame}[fragile]{Laboratorio en Python: Prueba $F$ e Intervalo de Confianza}
  \scriptsize
  Prueba de igualdad de varianzas, IC para $\sigma_1^2/\sigma_2^2$, y verificación de cobertura Monte Carlo:
  {\lstset{basicstyle=\fontsize{4pt}{4.8pt}\selectfont\ttfamily,aboveskip=0.1em,belowskip=0.1em}
  \lstinputlisting[firstline=39, lastline=67]{../../code/05_distribuciones_muestreo/05.05_fisher_f_distribution.py}}
\end{frame}

% Slide 10: Lab Python (Bloque 3)
\begin{frame}[fragile]{Laboratorio en Python: ANOVA Completo}
  \scriptsize
  Cálculo manual de ANOVA de una vía, verificado contra \texttt{scipy.stats.f\_oneway}:
  {\lstset{basicstyle=\fontsize{4pt}{4.8pt}\selectfont\ttfamily,aboveskip=0.1em,belowskip=0.1em}
  \lstinputlisting[firstline=70, lastline=96]{../../code/05_distribuciones_muestreo/05.05_fisher_f_distribution.py}}
\end{frame}

% Slide 11: Salida Terminal
\begin{frame}[fragile]{Laboratorio en Python: Salida en Terminal}
  \scriptsize
  Salida estándar generada al ejecutar el script del laboratorio en Python:
  \vspace{0.02cm}
  \begin{block}{Salida Estándar en Consola}
    \fontsize{4pt}{5pt}\selectfont
    \begin{verbatim}
=== Block 1: F-Distribution Properties ===
F_5,10: E(F)=1.2500 | Reciprocal KS p=0.643 | T^2~F_1,12 KS p=0.527

=== Block 2: F-Test and CI for sigma1^2/sigma2^2 ===
F=2.25, critical F_12,10,0.025=3.62 -> Reject H0? No
95% CI: [0.6214, 7.5905] | MC coverage: 0.9497 (nominal 0.95)

=== Block 3: One-Way ANOVA Verification ===
Means: A=12.0, B=9.0, C=15.0 | SS_treat=72.00, SS_error=16.00
F=20.2500, critical F_2,9,0.05=4.26 -> Reject H0? Yes
SciPy f_oneway: F=20.2500, p=0.000466
    \end{verbatim}
  \end{block}
\end{frame}

% Slide 12A: Ejercicio Nivel 1 (Enunciado)
\begin{frame}{Ejercicio en Clase --- Nivel 1: Fundamental (1/2)}
  \footnotesize
  \begin{block}{Problema 5.5.1 --- Media de $F_{5,10}$ (Enunciado)}
    Sea $F \sim F_{5,10}$. Calcule $E(F)$.
  \end{block}

  \vspace{0.15cm}
  \pause
  \begin{alertblock}{Planteamiento}
    \begin{itemize}\itemsep=0.04cm
      \item Use directamente $E(F)=d_2/(d_2-2)$.
    \end{itemize}
  \end{alertblock}
\end{frame}

% Slide 12B: Ejercicio Nivel 1 (Resolución)
\begin{frame}{Ejercicio en Clase --- Nivel 1: Fundamental (2/2)}
  \scriptsize
  Sustituimos $d_2=10$ directamente: \pause

  \vspace{0.05cm}
  \begin{block}{Cálculo}
    $$E(F) = \frac{10}{10-2} = \frac{10}{8} = 1.25.$$
  \end{block}

  \vspace{0.05cm}
  \pause
  \begin{alertblock}{Interpretación}
    \scriptsize
    Nótese que $E(F)$ depende únicamente de $d_2$ (el denominador), no de $d_1$. Cuando $d_2\to\infty$, $E(F)\to 1$, consistente con que el denominador $\chi^2_{d_2}/d_2 \to 1$.
  \end{alertblock}
\end{frame}

% Slide 13A: Ejercicio Nivel 2 (Enunciado)
\begin{frame}{Ejercicio en Clase --- Nivel 2: Operativo (1/2)}
  \footnotesize
  \begin{block}{Problema 5.5.4 --- Prueba $F$ de Varianzas (Enunciado)}
    Línea 1 ($n_1=13$): $S_1^2=45$. Línea 2 ($n_2=11$): $S_2^2=20$. Pruebe $H_0:\sigma_1^2=\sigma_2^2$ a nivel $\alpha=0.05$, usando $F_{12,10,0.025}\approx 3.62$.
  \end{block}

  \vspace{0.15cm}
  \pause
  \begin{alertblock}{Estrategia}
    \scriptsize
    Calcule $F=S_1^2/S_2^2$ y compárelo contra el valor crítico dado.
  \end{alertblock}
\end{frame}

% Slide 13B: Ejercicio Nivel 2 (Resolución)
\begin{frame}{Ejercicio en Clase --- Nivel 2: Operativo (2/2)}
  \scriptsize
  Calculamos el estadístico $F$: \pause

  \vspace{0.05cm}
  \begin{block}{Cálculo}
    $$F = \frac{45}{20} = 2.25$$
  \end{block}

  \vspace{0.05cm}
  \pause
  \begin{alertblock}{Decisión}
    \scriptsize
    Como $2.25 < 3.62$, \textbf{no se rechaza} $H_0:\sigma_1^2=\sigma_2^2$: no hay evidencia suficiente de que las varianzas de las dos líneas de producción difieran.
  \end{alertblock}
\end{frame}

% Slide 14A: Ejercicio Nivel 3 (Enunciado)
\begin{frame}{Ejercicio en Clase --- Nivel 3: Analítico (1/2)}
  \footnotesize
  \begin{block}{Problema 5.5.7 --- Identidad $T^2\sim F_{1,\nu}$ (Enunciado)}
    Sea $T\sim t_\nu$. Demuestre que $T^2\sim F_{1,\nu}$, partiendo de $T=Z/\sqrt{\chi^2_\nu/\nu}$ con $Z\sim N(0,1)$, $Z\perp\chi^2_\nu$.
  \end{block}

  \vspace{0.1cm}
  \pause
  \begin{alertblock}{Estrategia}
    \scriptsize
    Eleve $T$ al cuadrado y reconozca que $Z^2\sim\chi^2_1$ por la definición de la chi-cuadrada.
  \end{alertblock}
\end{frame}

% Slide 14B: Ejercicio Nivel 3 (Resolución)
\begin{frame}{Ejercicio en Clase --- Nivel 3: Analítico (2/2)}
  \scriptsize
  Elevamos al cuadrado: \pause

  \vspace{0.05cm}
  \begin{block}{Demostración}
    \begin{align*}
      T^2 = \frac{Z^2}{\chi^2_\nu/\nu} = \frac{\chi^2_1/1}{\chi^2_\nu/\nu}, \qquad \text{ya que } Z^2\sim\chi^2_1. \quad \qedhere
    \end{align*}
  \end{block}

  \vspace{0.05cm}
  \pause
  \begin{alertblock}{Interpretación}
    \scriptsize
    Esta identidad muestra que las pruebas $t$ y $F$ para un solo predictor son equivalentes: el cuadrado del estadístico $t$ en una regresión simple es exactamente el estadístico $F$ correspondiente, con $1$ grado de libertad en el numerador.
  \end{alertblock}
\end{frame}

% Slide 15A: Ejercicio Nivel 4 (Enunciado)
\begin{frame}{Ejercicio en Clase --- Nivel 4: Desafiante (1/2)}
  \footnotesize
  \begin{block}{Problema 5.5.9 --- ANOVA Completo (Enunciado)}
    Tres métodos de enseñanza, $n=4$ c/u: A: $10,12,14,12$; B: $8,9,11,8$; C: $15,16,14,15$. Calcule $\text{SC}_{\text{trat}}$, $\text{SC}_{\text{error}}$, $F$, y decida a $\alpha=0.05$ usando $F_{2,9,0.05}\approx 4.26$.
  \end{block}

  \vspace{0.15cm}
  \pause
  \begin{alertblock}{Estrategia}
    \scriptsize
    Calcule primero las medias de grupo y la media global, luego las sumas de cuadrados entre y dentro de grupos.
  \end{alertblock}
\end{frame}

% Slide 15B: Ejercicio Nivel 4 (Resolución)
\begin{frame}{Ejercicio en Clase --- Nivel 4: Desafiante (2/2)}
  \scriptsize
  Medias de grupo: $12$, $9$, $15$; media global: $12$. \pause

  \vspace{0.05cm}
  \begin{block}{Cálculo}
    \scriptsize
    \begin{align*}
      \text{SC}_{\text{trat}} &= 4(0)^2+4(-3)^2+4(3)^2 = 72, \qquad \text{SC}_{\text{error}} = 8+6+2=16.
    \end{align*}
    \begin{align*}
      F = \frac{72/2}{16/9} = \frac{36}{1.778} \approx 20.25.
    \end{align*}
  \end{block}

  \vspace{0.05cm}
  \pause
  \begin{alertblock}{Decisión}
    \scriptsize
    Como $20.25 > 4.26$, \textbf{se rechaza} $H_0:\mu_A=\mu_B=\mu_C$: hay fuerte evidencia de que al menos un método de enseñanza produce resultados diferentes.
  \end{alertblock}
\end{frame}

% Slide 16: Cierre
\begin{frame}{Cierre de Sección: la $F$ de Fisher-Snedecor}
  \begin{reflexion}
    Con la distribución $F$ completamos el trío $\chi^2$--$t$--$F$: todas surgen de combinar normales estándar independientes, y todas resuelven un problema práctico específico de la inferencia con muestras normales. Este es un cierre elegante: cada distribución del capítulo se construye a partir de las anteriores.
  \end{reflexion}

  \vspace{0.2cm}
  \pause
  \begin{block}{Lecciones Clave}
    \begin{itemize}\itemsep=0.08cm
      \item \textbf{Definición:} $F=(\chi^2_{d_1}/d_1)/(\chi^2_{d_2}/d_2)$; $E(F)=d_2/(d_2-2)$; $1/F_{d_1,d_2}\sim F_{d_2,d_1}$.
      \item \textbf{Prueba de varianzas:} $F=S_1^2/S_2^2\sim F_{n_1-1,n_2-1}$ bajo $H_0:\sigma_1^2=\sigma_2^2$.
      \item \textbf{ANOVA:} $F=\text{SC}_{\text{trat}}/(k-1) \div \text{SC}_{\text{error}}/(n-k)$ compara $k$ medias simultáneamente.
    \end{itemize}
  \end{block}
\end{frame}

% Slide 17: Perspectiva Modular
\begin{frame}{Perspectiva Modular --- Cierre del Capítulo 05}
  \scriptsize
  \begin{retoclase}
    Con la Sección 05.05 concluimos el Capítulo 05: Distribuciones de Muestreo, habiendo formalizado la distribución de $\bar X$ y $S^2$, el TLC asintótico, y las distribuciones $\chi^2$, $t$ y $F$. El Capítulo 06 --- Estimación y su Relación con Ciencia de Datos --- usará estas herramientas para construir estimadores puntuales (momentos, máxima verosimilitud) e intervalos de confianza generales, cerrando el puente entre la teoría de probabilidad y la inferencia estadística aplicada.
  \end{retoclase}

  \vspace{0.08cm}
  \pause
  \begin{alertblock}{Preparación Recomendada}
    \begin{itemize}\itemsep=0.06cm
      \item Repasa los conceptos de insesgadez y consistencia (Sección 05.01): son los criterios centrales para evaluar estimadores.
      \item Reflexiona sobre cómo los intervalos de confianza que construimos hoy y en secciones anteriores son casos particulares de un marco más general de estimación.
    \end{itemize}
  \end{alertblock}
\end{frame}

\end{document}
